%% toc: P3 26 26.5,26.6,26.7
\chapter{The Vaught Ordinal and the Countable-or-Continuum Dichotomy}
\label{sec:p03-w03-p03-016-vaught-ordinal}

Several formulations of the countable-model problem are distinguished in
\cref{def:p03-vaught-variants}: the first-order cardinal dichotomy, its
perfect-set strengthening, and the infinitary cardinal dichotomy for a
sentence \(\sigma\) of \(L_{\omega_1\omega}\) in a countable language.
\Cref{def:p03-vaught-scott-comparison} adds the Vaught ordinal
\(\operatorname{vo}(\sigma)\) and the \(\omega\)-Vaught conjecture.
This chapter relates the Vaught ordinal to the cardinal formulation. The input is a sentence
\(\sigma\); the conclusion is a cardinal dichotomy for its countable models
whenever \(\operatorname{vo}(\sigma)\) is countable. The chapter also
records what the continuum hypothesis does to the converse.

For an \(L_{\omega_1\omega}\)-sentence \(\sigma\) in a countable language,
write \(I(\sigma,\aleph_0)\) for the number of isomorphism types of countable
models of \(\sigma\), finite models included, exactly as \(I(T,\aleph_0)\)
is used for a first-order theory in \cref{def:p03-vaught-spectrum-notation}.
Models in the definition of \(\operatorname{vo}(\sigma)\) are countable
models. Satisfaction of \(L_{\omega_1\omega}\)-formulas is preserved by
isomorphisms, so isomorphic models realize the same
\(\Pi^{\mathrm{in}}_\beta\)-types.

\begin{hbtheorem}{thm:p03-w03-p03-016-vo-cardinal}{proved=annex:p03-w03-p03-016-vo-cardinal}{A countable Vaught ordinal gives the cardinal dichotomy}
Let \(\sigma\) be an \(L_{\omega_1\omega}\)-sentence in a countable language.
\begin{enumerate}
  \item If \(\operatorname{vo}(\sigma)<\omega_1\), then
        \(I(\sigma,\aleph_0)\leq\aleph_0\) or
        \(I(\sigma,\aleph_0)=2^{\aleph_0}\). More precisely, if the first
        alternative in the definition of \(\operatorname{vo}(\sigma)\) holds
        at some countable \(\beta\), then
        \(I(\sigma,\aleph_0)=2^{\aleph_0}\); if the second holds, then
        \(I(\sigma,\aleph_0)\leq\aleph_0\).
  \item If \(\sigma\) is a \(\Pi^{\mathrm{in}}_\alpha\)-sentence with
        \(\alpha<\omega_1\) and
        \(\operatorname{vo}(\sigma)\leq\alpha+\omega\), then
        \(\operatorname{vo}(\sigma)<\omega_1\). Hence the
        \(\omega\)-Vaught conjecture implies the infinitary cardinal
        dichotomy for every sentence.
  \item Assume the continuum hypothesis. Then every \(\sigma\) satisfies the
        cardinal dichotomy. Consequently the statement ``every
        \(\sigma\) satisfying the cardinal dichotomy has
        \(\operatorname{vo}(\sigma)<\omega_1\)'' is, under the continuum
        hypothesis, equivalent to the statement that every
        \(L_{\omega_1\omega}\)-sentence in a countable language has a
        countable Vaught ordinal.
\end{enumerate}
\end{hbtheorem}

Clause~3 is the precise sense in which a mere cardinal formulation is
weakened by the continuum hypothesis. The equivalence recorded by Gonzalez,
Rossegger and Turetsky between a countable Vaught ordinal and Vaught's
conjecture \cite[pp.~2, 6]{GonzalezRosseggerTuretsky2026Status} therefore
cannot be read, under the continuum hypothesis, with the bare cardinal
dichotomy in place of the conjecture unless every sentence has a countable
Vaught ordinal. Their introduction treats that general question as open
\cite[p.~1]{GonzalezRosseggerTuretsky2026Status}. Their argument for the
direction from the conjecture to smallness at every countable level
(Proposition~2.2) goes through Silver's theorem on Borel equivalence
relations rather than through a count of cardinals
\cite[p.~7]{GonzalezRosseggerTuretsky2026Status}. The handbook does not
prove that direction, and Theorem~\ref{thm:p03-w03-p03-016-vo-cardinal}
does not depend on it.

\begin{hbexample}{def:p03-w03-p03-016-independent-predicates}{Independent unary predicates have Vaught ordinal \(0\)}
Let the language consist of unary predicates \(P_0,P_1,\ldots\). For
disjoint finite sets \(A,B\subseteq\omega\) let \(\theta_{A,B}\) be
\[
 \exists x\Bigl(\bigwedge_{i\in A}P_i(x)\wedge
                 \bigwedge_{j\in B}\neg P_j(x)\Bigr),
\]
and let \(\sigma_P\) be the countable conjunction of all \(\theta_{A,B}\).
The structure whose domain is the set of finite subsets of \(\omega\), with
\(P_i\) interpreted as \(\{x:i\in x\}\), is a countable model. For every
\(S\subseteq\omega\) some countable model of \(\sigma_P\) has an element
lying in \(P_i\) exactly for \(i\in S\). These elements realize pairwise
different \(\Pi^{\mathrm{in}}_0\)-types, since each such type contains
\(P_i(x)\) or \(\neg P_i(x)\) for every \(i\). Hence continuum many
\(\Pi^{\mathrm{in}}_0\)-types are realized,
\(\operatorname{vo}(\sigma_P)=0\), and by clause~1 of
Theorem~\ref{thm:p03-w03-p03-016-vo-cardinal},
\(I(\sigma_P,\aleph_0)=2^{\aleph_0}\). The annex proof of that theorem
includes the verification of this example.
\end{hbexample}

\begin{hbwarning}{def:p03-w03-p03-016-vo-limits}{What the theorem does not supply}
Clause~1 is a cardinal conclusion. It does not produce a perfect set of
pairwise nonisomorphic models, so it does not decide the perfect-set
formulation of \cref{def:p03-vaught-variants}. Clause~2 uses only that
\(\alpha+\omega\) is countable; it says nothing about which sentences meet
the \(\omega\)-Vaught bound. The bound is known for the infinitary theories
with finitely many models \cite[p.~4, Corollary~4.10]{GonzalezRosseggerTuretsky2026Status},
as recorded in the dated status paragraph of
\cref{sec:p03-vaught-status-2026}. Neither clause, nor the never-two
theorem, settles the general conjecture.
\end{hbwarning}
